% !TeX root = ../Thesis.tex

%************************************************
\chapter{Discussion}\label{ch:discussion} % $\mathbb{ZNR}$
%************************************************
\section{Implications of Middleware-assisted Sharded Distributed Databases}
To address the three research questions introduced in chapter \ref{ch:introduction}, this thesis undertook the implementation and evaluation of a prototype middleware system that enabled sharding in databases, thereby creating a system of sharded distributed databases. The middleware's performance was assessed in terms of throughput, response time, and data distribution. This analysis identified several design-related limitations of the middleware. Despite these shortcomings and challenges, it delivered acceptable results, which leads to the conclusion that carefully designed sharded distributed databases have the potential to enhance the scalability of centralized databases. \\

\noindent From the load-test results presented in chapter \ref{evaluation}, the prototype of middleware designed for sharding exhibited acceptable performance in terms of system throughput. Upon further examination, it was determined that the middleware had limitations in managing multiple client requests due to its synchronous behavior. Each request had to wait for completion, occupying computing resources without utilizing them for other requests. During bulk read and write operations, the middleware experienced substantial high waiting times to receive results from the associated Redis databases. Consequently, the middleware's computing resources were engaged, but the process remained idle while waiting for responses from the Redis databases. As a single synchronous interface, the middleware was responsible for handling numerous client requests and orchestrating data management across multiple Redis databases. This also suggested the presence of a single bottleneck in the system, hindering optimal throughput. Despite these limitations, the middleware managed to achieve an average throughput of 80 Requests Per Second (RPS) when utilizing 6 Redis database instances, each containing 1 shard. This performance is considered acceptable considering the design goals of the middleware, as it was originally created as a single instance for communication and all operations were executed synchronously. There was no provision for the parallel execution of requests or concurrent computations to reduce the complexities of concurrency control and synchronization within the middleware. The primary objective was to implement a sharding approach using a middleware prototype to enable a functioning distributed database system, assess its performance, and identify challenges encountered during implementation. \\

\noindent Regarding response time, it was evident that the middleware connected to multiple Redis databases exhibited higher average response time compared to the baseline configuration of a single Redis database. When the middleware connected with multiple Redis databases was employed to synchronize data sharding and distribution, added network latencies arose due to the increased number of communication entities.The findings from the test execution indicated that the average response time observed for a single client connected to the middleware with six Redis database instance was approximately 9 milliseconds (ms) higher than that of the baseline, which consisted of a single client connected to one Redis database. This difference is reasonable given that the baseline involved only two communicating entities, while the middleware-assisted sharded distributed database system encompassed three and resulted in increased network latencies. When additional concurrent clients were connected, the middleware's average response time also exhibited a linear increase, which was an expected behavior due to the overloading of a single instance of middleware for handling multiple concurrent requests. \\

\noindent In terms of data distribution, the consistent hashing approach employing virtual nodes demonstrated effective data distribution. Test and evaluations in chapter \ref{evaluation}, which involved running the middleware with varying numbers of Redis databases and shards, indicated that the middleware successfully distributed 1000 key-value pairs across six Redis databases, achieving near-uniform data distribution without the risk of data hotspots. The evaluation showed a maximum of 183 key-value pairs assigned to one database and a minimum of 154 assigned to another, resulting in a maximum difference of merely 29 key-value pairs. \\

\noindent Considering the aforementioned findings and results, it can be concluded that the system of middleware-assisted sharded distributed databases leads to effective data distribution, which addresses the issue of overloaded single centralized databases. The middleware system adapts proficiently to multiple database node configurations, with each node containing a single shard. However, the evaluations of this thesis suggest that sharding introduces several challenges, including processing overhead in data and request distribution, compilation of partial results, and synchronization overhead across multiple database servers. Thus, sharded distributed databases should be suitably designed and optimized for peak performance, encompassing both throughput and response time.



%************************************************
\section{Limitations and Future Work}
